\section{Threats to validity}
\label{sec:methodology:threats}

Construct validity concerns the gap between the quantity an instrument records and the quantity the experiment intends to measure, and two outcomes in this design are exposed to it. Network transfer is recorded at the host boundary where each configuration issues its reads and collects its results (Section~\ref{sec:REPLACE-ME-primary-metrics}), so the figure reported is the volume the host observes rather than the bytes that traverse the wire. Transport-layer framing, encryption, and any retransmission sit below that boundary and are not separated out, so the transfer metric approximates the construct of data moved over the network rather than measuring it exactly. The same boundary applies uniformly to every configuration, so the approximation bears on an absolute figure more than on a comparison between configurations measured the same way. The k-nearest-neighbor workload raises a second concern, because distance ordering is served by the spatial index only when a bounding-box predicate accompanies it (Section~\ref{sec:REPLACE-ME-full-db-engines}), so an engine may rank candidates by a bounding-box distance proxy before refining to exact distance, and the proxy an engine applies is not guaranteed identical across configurations. Holding the query point and the returned count fixed bounds this divergence but does not remove it, since two engines can return the same cardinality while having ordered their candidates through different internal distance computations.

Internal validity asks whether an observed difference can be attributed to the factor under test rather than to a confound, and the dominant confound here is the performance variability of shared cloud infrastructure. Production cloud services vary over time and between nominally identical allocations, with run-to-run noise that is uncorrelated with the format-and-engine pairing under study \parencite{iosup2011_performance_variability_cloud_services}, and instances provisioned to the same specification have been shown to differ measurably in processor, input/output, and network throughput \parencite{schad2010_runtime_measurements_cloud}. Four measures established in Section~\ref{sec:REPLACE-ME-replication} limit this confound rather than assume it away. Configurations compared on the same workload are launched concurrently on identically provisioned instances, so the two members of a comparison share the same short-term conditions; the execution order of experiments is randomized under a fixed seed; warmup iterations bring each engine to steady operation before any timed measurement; and the distributed configuration provisions a fresh cluster for every run. Running test and control on one instance in randomized order is what allows a small difference to be read as an effect of the factor rather than as drift between separately timed runs.

External validity concerns how far the findings extend beyond the conditions under which they were produced, and three features of the design bound that reach. The experiment runs in a single cloud region, with a single provider, on one dataset domain of Norwegian building geometry. The orderings that follow from architecture are expected to carry beyond these conditions: which access pattern moves fewer bytes, and whether a distributed engine repays its coordination cost at a given scale, are consequences of how a format is laid out and how an engine reads it (Section~\ref{sec:REPLACE-ME-format-summary}; Section~\ref{sec:REPLACE-ME-compute-models}) rather than of the platform that hosts the run. Absolute wall-clock times and monetary costs do not carry in the same way, since they depend on provider-specific and region-specific rates, and public cloud providers differ in how predictable that performance is \parencite{leitner2016_patterns_chaos_iaas}. Generalization across data is the most limited axis, because building polygons from one national dataset need not represent the geometry complexity or spatial distribution of other inputs, and a workload whose cost is governed by spatial skew (Section~\ref{sec:REPLACE-ME-spatial-skew}), such as the distributed join, is the most sensitive to that difference.

Three asymmetries between configurations are accepted by the design rather than corrected, and each is recorded here so that it can be read alongside the comparison it affects. The Shapefile configuration is read from local disk rather than from object storage (Figure~\ref{fig:REPLACE-ME-design-cells}), because the format has no cloud-native read path, and this favors it on transfer relative to the configurations that fetch from object storage, so a transfer comparison involving Shapefile is read as a lower bound on the gap rather than as a like-for-like figure. The distributed configuration alone may move data across a region boundary and can therefore incur an egress charge that the single-node configurations, reading within one region, do not (Section~\ref{sec:REPLACE-ME-cost-model}). Warmup is omitted for the distributed configuration because each of its runs provisions a fresh cluster, and priming would multiply provisioning cost without leaving a persistent engine to prime, so its first-run variance is carried into the analysis rather than removed by warmup (Section~\ref{sec:REPLACE-ME-replication}). None of the three is hidden inside an aggregate; each is reported against the workload and dimension it touches, so that a reader can discount the affected comparison rather than take it at face value.